\section{Diagonalization --- Not Time Constructible}
%%% generated by the blueprint pipeline from TCSlib.Complexity.Diagonalization.NotTimeConstructible
\subsection*{Diagonalization --- Not Time Constructible: Introduction}

Declarations of the Lean module \texttt{TCSlib.Complexity.Diagonalization.NotTimeConstructible}.

\subsection*{Diagonalization --- Not Time Constructible: Results}

\begin{definition}[Rank of a Boolean string]
\lean{Complexity.ntcRank}
\label{Complexity.ntcRank}
\leanok
\uses{Turing.codeBitsNat}
For a finite Boolean string $s$, define its rank to be the natural number represented by
the bit code of $s$ with a true bit appended, minus one, where subtraction is truncated
at zero.
\end{definition}

\begin{definition}[Binary string associated to a natural number]
\lean{Complexity.ntcStr}
\label{Complexity.ntcStr}
\leanok
For a natural number $n$, define its associated Boolean string by taking the binary
digits of $n+1$ and removing the final digit.
\end{definition}

\begin{definition}[Halting-adjusted natural number]
\lean{Complexity.ntcT}
\label{Complexity.ntcT}
\leanok
\uses{Complexity.HALT, Complexity.TimeHierarchy.code, Complexity.ntcStr}
For a natural number $n$, define a natural number equal to $n+1$ if the time-hierarchy
machine halts on the Boolean string associated to $n$, and equal to $n$ otherwise.
\end{definition}

\begin{theorem}[Bit Encoding and Base-Two Digits]
\lean{Complexity.ntc_codeBitsNat_eq_ofDigits}
\label{Complexity.ntc_codeBitsNat_eq_ofDigits}
\leanok
\uses{Turing.codeBitsNat}
\difficulty{4}
For every finite bit string $w$, its least-significant-bit-first binary value equals the
base-two numeral whose digit list is obtained by replacing each false bit by $0$ and
each true bit by $1$.
\end{theorem}

\begin{theorem}[Appending a One in Binary Encoding]
\lean{Complexity.ntc_bits_codeBitsNat_append_true}
\label{Complexity.ntc_bits_codeBitsNat_append_true}
\leanok
\uses{Complexity.ntc_codeBitsNat_eq_ofDigits, Turing.codeBitsNat}
\difficulty{5}
For every finite bit string $s$, the binary expansion of the natural number encoded by
appending $1$ to $s$, with bits interpreted in least-significant-bit-first order, is
exactly $s$ followed by $1$.
\end{theorem}

\begin{theorem}[Binary Expansion of the Successor Rank]
\lean{Complexity.ntc_bits_rank_succ}
\label{Complexity.ntc_bits_rank_succ}
\leanok
\uses{Complexity.ntcRank, Complexity.ntc_bits_codeBitsNat_append_true, Turing.codeBitsNat}
\difficulty{4}
For every finite bit string $s$, let its rank be the natural number encoded by appending
$1$ to $s$ and then subtracting $1$. The binary expansion of one more than this rank is
exactly $s$ followed by $1$.
\end{theorem}

\begin{theorem}[String Recovery from Rank]
\lean{Complexity.ntc_str_rank}
\label{Complexity.ntc_str_rank}
\leanok
\uses{Complexity.ntcRank, Complexity.ntcStr, Complexity.ntc_bits_rank_succ}
\difficulty{3}
For every bit string $s$, form the binary expansion of $r(s)+1$, with bits listed
least-significant first, and remove its last bit. The resulting bit string is exactly
$s$.
\end{theorem}

\begin{theorem}[Parity of the modified time bound]
\lean{Complexity.ntc_bodd_T}
\label{Complexity.ntc_bodd_T}
\leanok
\uses{Complexity.HALT, Complexity.TimeHierarchy.code, Complexity.ntcStr, Complexity.ntcT}
\difficulty{3}
For every natural number $n$, the parity of $n$ plus $1$ if the halting predicate for
the fixed time-hierarchy code accepts the binary string associated to $n$, and plus $0$
otherwise, equals the exclusive-or of the parity of $n$ and the truth value of that
predicate.
\end{theorem}

\begin{theorem}[Primitive Recursiveness of Rank]
\lean{Complexity.ntc_primrec_rank}
\label{Complexity.ntc_primrec_rank}
\leanok
\uses{Complexity.ntcRank, Turing.codePrimBitsNat}
\difficulty{3}
For a bit string $s$, let $r(s)$ be the natural number obtained by appending a $1$ to
$s$, interpreting the resulting string in least-significant-bit-first binary, and
subtracting $1$. The function $s\mapsto r(s)$ is primitive recursive.
\end{theorem}

\begin{theorem}[Primitive Recursive Unary Output]
\lean{Complexity.ntc_primrec_emit}
\label{Complexity.ntc_primrec_emit}
\leanok
\uses{Complexity.ntcRank, Complexity.ntc_primrec_rank}
\difficulty{3}
The function sending each finite bit string $s$ to the list consisting of
$\operatorname{rank}(s)$ copies of $1$ is primitive recursive, where
$\operatorname{rank}(s)$ is the least-significant-bit-first binary value of $s$ with a
final $1$ appended, minus $1$.
\end{theorem}

\begin{theorem}[Primitive Recursiveness of Rank Parity]
\lean{Complexity.ntc_primrec_parity}
\label{Complexity.ntc_primrec_parity}
\leanok
\uses{Complexity.ntcRank, Complexity.ntc_primrec_rank}
\difficulty{3}
For a bit string $s$, let $r(s)$ be the natural number obtained by appending a $1$ to
$s$, interpreting the resulting string in least-significant-bit-first binary, and
subtracting $1$. The function sending $s$ to the singleton list containing whether
$r(s)$ is odd is primitive recursive.
\end{theorem}

\begin{theorem}[Primitive Recursiveness of the List Head]
\lean{Complexity.ntc_primrec_head}
\label{Complexity.ntc_primrec_head}
\leanok
\difficulty{3}
The function that sends a Boolean list $w$ to the singleton list containing its first
bit, taking the bit to be false when $w$ is empty, is primitive recursive.
\end{theorem}

\begin{theorem}[Primitive Recursiveness of the Negated List Head]
\lean{Complexity.ntc_primrec_nhead}
\label{Complexity.ntc_primrec_nhead}
\leanok
\difficulty{3}
The function that sends a Boolean list $w$ to the singleton list containing the negation
of its first bit, taking the first bit to be false when $w$ is empty, is primitive
recursive.
\end{theorem}

\begin{theorem}[Failure of time constructibility]
\lean{Complexity.ntc_not_timeConstructible}
\label{Complexity.ntc_not_timeConstructible}
\leanok
\uses{Complexity.HALT, Complexity.HALT_not_computable, Complexity.TimeConstructible, Complexity.TimeHierarchy.code, Complexity.ntcRank, Complexity.ntcT, Complexity.ntc_bodd_T, Complexity.ntc_primrec_emit, Complexity.ntc_primrec_head, Complexity.ntc_primrec_nhead, Complexity.ntc_primrec_parity, Complexity.ntc_str_rank, Turing.FinTM.ComputesFunInTime.computes, Turing.FinTM.ComputesInTime, Turing.FinTM.exists_comp_partial, Turing.FinTM.exists_cond, Turing.codePrim_machine}
\difficulty{6}
Define $T(n)=n+1$ if the halting predicate for the fixed time-hierarchy code accepts the
binary string associated to $n$, and define $T(n)=n$ otherwise. Then $T$ is not
time-constructible: it is false that both $n\le T(n)$ for every natural number $n$ and
that there exist a positive natural number $c$ and a finite-state Turing machine $M$
over the binary alphabet such that, for every finite binary input $x$, $M$ halts with
output the binary expansion of $T(|x|)$ within $c(T(|x|)+1)$ steps.
\end{theorem}

\begin{theorem}[A non-time-constructible dominating function]
\lean{Complexity.exists_not_timeConstructible}
\label{Complexity.exists_not_timeConstructible}
\leanok
\uses{Complexity.TimeConstructible, Complexity.ntcT, Complexity.ntc_not_timeConstructible}
\difficulty{1}
There exists a function $T:\mathbb{N}\to\mathbb{N}$ such that $n\le T(n)$ for every $n$,
but $T$ is not time constructible. Here time constructibility means both that $T(n)\ge
n$ and that some finite binary Turing machine, on every input of length $n$, outputs the
binary expansion of $T(n)$ and halts within $c(T(n)+1)$ steps for some positive integer
$c$.
\end{theorem}
